\documentclass[11pt]{article}
\usepackage[T1]{fontenc}
\usepackage[utf8]{inputenc}
\usepackage{amsmath,amssymb,amsthm}
\usepackage[margin=1in]{geometry}
\usepackage{hyperref}
\newtheorem{theorem}{Theorem}
\newtheorem{lemma}[theorem]{Lemma}
\newtheorem{corollary}[theorem]{Corollary}
\newcommand{\OO}{\mathcal O}
\newcommand{\GG}{\mathcal G}
\newcommand{\FF}{\mathcal F}
\title{Relative regularity over an affine base and a Frobenius generation bridge}
\author{Unreviewed mathematical working paper}
\date{8 October 2026}
\begin{document}
\maketitle
\begin{abstract}
We prove the zero-regularity generation criterion on projective space over
a finite type algebra over an arbitrary field. The argument includes
finite ground fields and does not require flatness of the coherent sheaf
over the base. We then obtain uniform-in-index generation for the
Frobenius source sheaves occurring in a proposed singular test-ideal
approximation. The latter application assumes its exact trace presentation;
it proves no test-ideal identity, subadditivity, or envelope theorem.
The complete original problem P12.3 remains open.
\end{abstract}

\section{Statement and foundations}
Let $k$ be any field, let $R$ be a finite type $k$-algebra, and put
$P=\mathbf P^N_R$. Twisting a coherent sheaf on $P$ means tensoring
with $\OO_P(t)$. All cohomology and global sections are on the entire
scheme, with its affine base $\operatorname{Spec}R$; they are not fibrewise
cohomology groups. A sheaf is globally generated if its evaluation map
$H^0(P,\GG)\otimes_R\OO_P\longrightarrow\GG$ is surjective.

\begin{theorem}\label{thm:relative}
If a coherent sheaf $\GG$ on $\mathbf P^N_R$ satisfies
\[
 H^i(P,\GG(-i))=0\qquad(1\leq i\leq N),
\]
then $\GG$ is globally generated. More precisely, for every integer
$t\geq0$,
\[
 H^i(P,\GG(t-i))=0\quad(i>0),
\]
and the multiplication map
\[
 H^0(P,\GG(t))\otimes_R R^{N+1}\longrightarrow
 H^0(P,\GG(t+1))
\]
defined by the homogeneous coordinates is surjective.
\end{theorem}

The foundations used below are coherent sheaf cohomology on a separated
Noetherian scheme, affine acyclicity and its finite affine Cech complex;
finite associated primes for a finite module over a Noetherian ring and
the zero-divisor criterion using those primes; and the following precise
projective-space result. Stacks, Lemma 30.14.1, Tag
\href{https://stacks.math.columbia.edu/tag/01YS}{01YS}, states for
Noetherian $R$ that every coherent $\GG$ on $\mathbf P^N_R$ is a quotient
of a finite direct sum of $\OO_P(d_j)$, that its cohomology vanishes
above degree $N$ and is finite over $R$, and that positive cohomology of
high twists vanishes. In particular high twists are globally generated:
twist the displayed surjection until all its source line bundles are
generated by homogeneous monomials. These exact facts hold over $R$,
with no flatness or field-base assumption. The immediate proof of Tag
01YS has been read in full; its deeper foundational references are not
claimed to have been freshly reconstructed in their entirety.
For the scheme-theoretic associated-point step we use the complete
statements and immediate proofs of Stacks 31.2.2, 31.2.5, and 31.2.10
in Tag \href{https://stacks.math.columbia.edu/tag/02OI}{02OI}:
associated points agree with associated primes on Noetherian affine
charts, are finite on quasi-compact opens for a coherent sheaf, and
a morphism of quasi-coherent sheaves injective at all associated
points of its source is injective everywhere. These units have been
freshly read; their recursive algebra foundations are cited.

\section{Flat extension and a regular hyperplane}
For any extension field $k'/k$ let $R'=R\otimes_k k'$ and let
$P'=\mathbf P^N_{R'}$. These are Noetherian because $R'$ is finite type
over $k'$. The map $R\to R'$ is faithfully flat. On the standard affine
cover of $P$, all intersections are affine. The Cech terms of a
quasi-coherent sheaf commute with flat base extension. Tensoring that
finite complex by the flat $R$-module $R'$ commutes with its kernels,
images, and cohomology. Consequently, for every $i$ and twist $t$,
\[
 H^i(P,\GG(t))\otimes_R R'
   =H^i(P',\GG'(t)).                                      \tag{1}
\]
If $\GG'$ is globally generated, (1) for $i=0$ identifies its evaluation
map with the pullback of the evaluation map for $\GG$. Its cokernel
therefore vanishes after faithfully flat pullback, hence vanishes.
This proves descent of global generation without a fibrewise base-change
theorem. We may replace $k$ by an infinite extension field to prove
Theorem~\ref{thm:relative}; this replacement does not assert any
Frobenius or test-ideal base-change identity.

Suppose now $k$ is infinite and $N\geq1$. The associated points of
$\GG$ form a finite set. For each such point $x$, the set of vectors
$(c_0,\ldots,c_N)\in k^{N+1}$ for which
\[
 l=\sum_{j=0}^N c_j X_j
\]
vanishes in the residue fibre of $\OO_P(1)$ at $x$ is a proper
$k$-linear subspace: at least one coordinate is nonzero there. Finitely
many proper linear subspaces cannot cover $k^{N+1}$. Choose a nonzero
vector outside their union. At every associated point $l$ is a unit
after trivializing $\OO_P(1)$, so multiplication is injective there.
Stacks 31.2.10 then makes multiplication by $l$ injective on
$\GG(-1)$ everywhere. One nonzero $c_j$ is a unit of $R$;
a constant invertible linear coordinate change makes $l$ a coordinate.
Thus its hyperplane $H$ is $\mathbf P^{N-1}_R$. With
$i:H\hookrightarrow P$ and $\FF=i^*\GG$, we have, for every integer $d$,
\[
 0\longrightarrow\GG(d-1)\xrightarrow{l}\GG(d)
   \longrightarrow i_*\FF(d)\longrightarrow0.              \tag{2}
\]
Neither $\GG$ nor $\FF$ is required to be $R$-flat.

\section{Complete induction proof}
We prove all assertions of Theorem~\ref{thm:relative} simultaneously
by induction on $N$. For $N=0$, projective space is $\operatorname{Spec}R$,
its twists are trivial, higher cohomology vanishes, and the single
coordinate gives the asserted multiplication map.

For $N>0$, use (1) to pass to an infinite field and construct (2).
For $1\leq i\leq N-1$, its cohomology sequence contains
\[
 H^i(P,\GG(-i))\longrightarrow H^i(H,\FF(-i))
  \longrightarrow H^{i+1}(P,\GG(-i-1)).
\]
The two outside groups vanish by the stated zero-regularity assumptions.
Hence $\FF$ is zero-regular on $H$. The induction hypothesis gives
all the higher-twist vanishings and the multiplication maps for $\FF$.

Next prove by induction on $t\geq0$ that
$H^i(P,\GG(t-i))=0$ for every $i>0$. For $t=0$ this is the hypothesis
for $i\leq N$ and projective-space cohomological dimension for $i>N$.
For the induction step, (2) gives the exact segment
\[
 H^i(P,\GG(t-1-i))\longrightarrow H^i(P,\GG(t-i))
  \longrightarrow H^i(H,\FF(t-i)).
\]
Its first group vanishes by the preceding $t$-induction and its third
by the $N$-induction. This proves the desired vanishing. In particular
$H^1(P,\GG(t-1))=0$ for all $t\geq0$, so restriction
$H^0(P,\GG(t))\to H^0(H,\FF(t))$ is surjective.

Let $a\in H^0(P,\GG(t+1))$. Its restriction to $H$ is, by induction,
a sum of products of sections of $\FF(t)$ and the $N$ remaining
coordinate sections on $H$. Lift those sheaf sections to $\GG(t)$
using the just-proved surjectivity, and lift the coordinate sections
by the fixed coordinate system on $P$. Subtracting these products
from $a$ leaves a section restricting to zero on $H$. Exactness of
(2) on global sections makes this remaining section $l b$ for some
$b\in H^0(P,\GG(t))$. This proves the multiplication map for $\GG$.
The description is invariant under the constant coordinate change.

Iterating, every section of $\GG(d)$, $d\geq0$, is a sum of homogeneous
degree-$d$ polynomials times sections of $\GG$. Choose $d$ large enough
that $\GG(d)$ is globally generated, as justified from Tag 01YS.
Let $\mathcal I$ be the image subsheaf of evaluation for $\GG$.
The multiplication assertion makes all global sections of $\GG(d)$
belong to $\mathcal I(d)$. Since those sections generate $\GG(d)$,
$\mathcal I(d)=\GG(d)$, hence $\mathcal I=\GG$. This proves global
generation over the infinite field, and its descent was proved in (1).
The theorem follows.

\begin{corollary}\label{cor:Y}
Let $Y\hookrightarrow\mathbf P^N_R$ be a closed immersion and let
$A=\OO_P(1)|_Y$. If a coherent sheaf $Q$ on $Y$ satisfies
$H^i(Y,Q\otimes A^{-i})=0$ for $1\leq i\leq N$, then $Q$ is globally
generated on $Y$.
\end{corollary}
\begin{proof}
The closed immersion is affine, its coherent pushforward is coherent,
and the affine Cech complexes give
$H^i(P,i_*Q(-i))=H^i(Y,Q\otimes A^{-i})$. Apply the theorem to $i_*Q$.
Pull back its surjective evaluation map to $Y$, and use
$i^*i_*Q=Q$. These statements can also be checked on affine charts
as restriction of a finite module along a quotient ring.
\end{proof}

\section{Frobenius source sheaves with a fixed generation twist}
Now assume $k$ is perfect of characteristic $p>0$, $R$ is finite type
over $k$, and $Y$ is projective over $R$. Choose a relatively very ample
line bundle $A=\OO_Y(H)$ and a closed immersion using its finitely
many generating sections into $\mathbf P^N_R$. Absolute Frobenius of
$Y$ is finite: this is the usual finite-presentation consequence of
perfectness and finite type. Let $L$ be a Cartier divisor and let
$M_0,\ldots,M_{b-1}$ be a fixed finite list of coherent sheaves on $Y$.
Choose an integer $c_0$ such that
\[
 (c_0-i)H-L\text{ is ample}\qquad(1\leq i\leq N).           \tag{3}
\]
Such a choice exists by the ample twisting property; it depends on
$H,L,N$ and not on the finite coherent list. For example first make
$\OO_Y(\nu H-L)$ globally generated, then its tensor product with
the very ample $A$ is ample, and take $c_0\geq\nu+N+1$.

\begin{corollary}\label{cor:Frob}
For all sufficiently large $e$ and every $j\in\{0,\ldots,b-1\}$,
the coherent sheaf
\[
 Q_{e,j}=F_*^e\bigl(M_j\otimes
                \OO_Y(p^e(c_0H-L))\bigr)
\]
is globally generated. The large-$e$ threshold may depend on the
$M_j$, whereas $c_0$ need not.
\end{corollary}
\begin{proof}
Write $q=p^e$. Projection formula for the finite Frobenius morphism gives
\[
 Q_{e,j}\otimes A^{-i}
 =F_*^e\bigl(M_j\otimes\OO_Y(q((c_0-i)H-L))\bigr).
\]
The topological map of absolute Frobenius is the identity; equivalently
its affine pushforward is cohomologically exact on quasi-coherent
sheaves. Hence the cohomology on the left is, as an abelian group,
the cohomology of the sheaf inside the pushforward. Relative Serre
vanishing over the Noetherian affine base applies to the coherent $M_j$
and the ample line bundle in (3). There are only finitely many pairs
$(i,j)$; choose their common Serre threshold and then $p^e$ above it.
All vanishings in Corollary~\ref{cor:Y} hold. Its conclusion proves
the assertion. No local freeness or Cohen--Macaulay property of $M_j$
is used.
\end{proof}

The Serre theorem invoked in this last proof has its exact standard
meaning: for a proper scheme over a Noetherian affine base and an ample
line bundle $B$, every coherent $M$ has
$H^i(Y,M\otimes B^q)=0$ for all $i>0$ and all sufficiently large
integers $q$. It is a published foundation, not a new statement about
test ideals; see Stacks, Lemma 30.17.1, Tag
\href{https://stacks.math.columbia.edu/tag/0B5U}{0B5U}, and Lemma
30.16.1, Tag \href{https://stacks.math.columbia.edu/tag/0B5T}{0B5T}.
The complete statements and immediate proofs of both have been freshly
read. The latter proof takes a very ample power, pushes to projective
space, and uses finitely many residue classes of the exponent together
with Tag 01YS. Its deeper ample-embedding and cohomology foundations
are not falsely described as newly proved here. No projective-over-a-field theorem is
substituted for this affine-base hypothesis.

\section{The exact conditional trace application and remaining scope}
Suppose a normal finite type $Y$ has a Weil canonical divisor $K_Y$,
a Cartier divisor $C$, and Cartier $L$. Suppose for a fixed positive
integer $m$ there exist an integer $\ell\geq1$ and a Cartier $D_0$
such that, for every integer threshold $E$, the sheaf
$\mathcal T_m\otimes\OO_Y(mC+c_0H)$ is a quotient of a finite sum of
the following sheaves with possibly different exponents $e\geq E$:
\[
 F_*^e\OO_Y\bigl(q(c_0H-L)+K_Y-D_0-j_eC\bigr),\qquad
 j_e=\ell\left\lceil\frac{qm}{\ell}\right\rceil-qm
       \in\{0,\ldots,\ell-1\}.                            \tag{4}
\]
Here $\OO_Y$ of a Weil divisor is the coherent divisorial sheaf.
Use Corollary~\ref{cor:Frob} with the finite list
$M_j=\OO_Y(K_Y-D_0-jC)$. Choose $E$ beyond the finite-list Serre
threshold, and take the assumed finite quotient using only exponents
$e\geq E$. Every summand is then globally generated, which proves
global generation of
the displayed target. A quotient of globally generated sheaves is
globally generated by the composed evaluations. Allowing $D_0$ and
$\ell$ to depend on $m$ changes this threshold but not $c_0$.
Thus a genuinely exact trace presentation of type (4) yields the
uniform-in-$m$ generation twist.

This is a conditional implication with an explicit antecedent. We have
not proved that $\mathcal T_m$ is an asymptotic test ideal, that a
scheme-global test divisor gives (4), that every arbitrarily late tail
of exponents gives the required total image, or that the indicated
quotient maps exist. A tail-sum presentation is sufficient; a single
exponent need not produce the entire target.
These are separate gates in working paper377. In that application $H$
can be chosen very ample by taking a sufficiently high power in its
primary-model construction. If desired, increasing its $c_0$ to make
(3) hold through the embedding dimension $N$ supplies the hypotheses
here without invoking cohomological-dimension reduction from $N$ to
$\dim Y$. This paper does not modify the already frozen working
paper377 or send feedback to its active independent review.

There is no use of a smoothness assumption, canonical invertibility,
uncorrected singular subadditivity, or a new Picard theorem in the
regularity result. The Frobenius corollary does require the stated
perfect-characteristic-$p$ finite type category. General original
P12.3 envelopes, Jacobian and boundary estimates, all-support positivity
recognition, mixed characteristic, and nonalgebraic excellent scope
receive no solution or correctness grant from this lemma.
No original-problem count or accepted-paper count is increased.
\end{document}
